\chapter*{How to read this script}
\addcontentsline{toc}{chapter}{How to read this script}
\markboth{How to read this script}{How to read this script}

This script is for someone who knows how to program but is new to medical imaging, to Kaggle competitions,
or to both. It follows one competition from start to finish: what the data looks like, which ideas we tried
and why, what each one scored, and what the team that finished first did.

\textbf{You do not need to read it in order.} Every chapter starts with a one-minute summary and ends with a
checkpoint. The pictures carry most of the story.

\section*{The boxes}

The text is broken up by coloured boxes. Each box has an icon and a word, so the meaning survives a
black-and-white printer.

\begin{tldr}
Three to five bullets at the start of every chapter. If you only read these, you get the whole story.
\end{tldr}
\begin{keyidea}
The one sentence to remember from a section.
\end{keyidea}
\begin{why}
The reason behind a choice. Read these when something feels arbitrary.
\end{why}
\begin{wordcheck}
A new word, explained in plain language the first time it appears. All of them are collected in the glossary
(Appendix~\ref{app:glossary}).
\end{wordcheck}
\begin{watchout}
A mistake that is easy to make, often one we made ourselves.
\end{watchout}
\begin{innumbers}
The figures that matter, in one place.
\end{innumbers}
\begin{example}
A small worked example, usually about one imaginary knee called \textbf{Study K}.
\end{example}
\begin{checkpoint}
\begin{itemize}
\done A short list of what you should now be able to explain.
\quiz Two or three questions to test yourself. The answers are in Appendix~\ref{app:answers}.
\end{itemize}
\end{checkpoint}

\section*{Three ways through}

\begin{center}
\begin{tikzpicture}[x=0.9mm,y=1mm,
  route/.style={draw=#1, line width=0.8pt, rounded corners=2mm, fill=white, minimum height=10mm,
                text width=34mm, align=center, font=\small},
  lab/.style={font=\small\bfseries, anchor=east}]
  \node[lab] at (-2,20) {\faStopwatch\ 15 minutes};
  \node[route=tldrline] at (24,20) {every \textbf{In one minute} box};
  \node[route=tldrline] at (70,20) {every \textbf{Checkpoint}};
  \node[route=tldrline] at (116,20) {Chapter~\ref{ch:usvsthem}: us vs.\ them};
  \node[lab] at (-2,6) {\faCoffee\ one hour};
  \node[route=keyline] at (24,6) {Part~I: the problem};
  \node[route=keyline] at (70,6) {Chapter~\ref{ch:experiments}: our experiments};
  \node[route=keyline] at (116,6) {Part~IV: the winners};
  \node[lab] at (-2,-8) {\faBook\ properly};
  \node[route=wordline, text width=118mm] at (70,-8) {front to back; Part~II explains the machine learning that Parts~III and IV use};
  \foreach \y in {20,6} {\draw[-{Stealth}, muted] (46.5,\y) -- (47.5,\y); \draw[-{Stealth}, muted] (92.5,\y) -- (93.5,\y);}
\end{tikzpicture}
\end{center}

\section*{The whole story on one page}

\begin{figure}[H]
\centering
\begin{tikzpicture}[x=1mm,y=1mm,
  stage/.style={draw=grid, fill=white, rounded corners=1.5mm, text width=29mm, minimum height=21mm,
                align=center, font=\footnotesize, inner sep=2pt},
  head/.style={font=\footnotesize\sffamily\bfseries, text=muted},
  ar/.style={-{Stealth[length=2mm]}, line width=0.8pt}]
  % column heads
  \foreach \x/\t in {0/{\faFile\ 1. Labels}, 33/{\faLayerGroup\ 2. Images in}, 66/{\faBrain\ 3. Model}, 99/{\faTrophy\ 4. Score}}
    \node[head] at (\x,17) {\t};
  % our lane
  \node[font=\small\bfseries, text=us, rotate=90] at (-21,0) {Us};
  \node[stage, draw=us, line width=0.8pt] (u1) at (0,0) {an open model reads the \textbf{report text}\\[1pt] agrees with doctors\\ \textbf{0.879} AUC};
  \node[stage, draw=us, line width=0.8pt] (u2) at (33,0) {3 series,\\ 24 slices each,\\ 256\,px, whole knee};
  \node[stage, draw=us, line width=0.8pt] (u3) at (66,0) {small CNNs\\ (EfficientNet-B0,\\ ConvNeXt-Nano)};
  \node[stage, draw=us, line width=0.8pt] (u4) at (99,0) {own model 0.918 public\\ best entry \textbf{0.927} private\\ rank 2,968 / 6,425};
  % winners lane
  \node[font=\small\bfseries, text=them, rotate=90] at (-21,-26) {Winners};
  \node[stage, draw=them, line width=0.8pt] (w1) at (0,-26) {AI agents read the \textbf{report and the images}\\[1pt] \textbf{0.94} AUC};
  \node[stage, draw=them, line width=0.8pt] (w2) at (33,-26) {5 series, all slices,\\ 384--512\,px,\\ + \textbf{anatomy boxes}};
  \node[stage, draw=them, line width=0.8pt] (w3) at (66,-26) {larger CNNs that\\ \textbf{look inside the boxes}};
  \node[stage, draw=them, line width=0.8pt] (w4) at (99,-26) {\textbf{0.960} private\\ rank 1\\ (one model: 0.954)};
  \foreach \a/\b in {u1/u2,u2/u3,u3/u4} \draw[ar, us] (\a) -- (\b);
  \foreach \a/\b in {w1/w2,w2/w3,w3/w4} \draw[ar, them] (\a) -- (\b);
  % the biggest gap
  \draw[decorate, decoration={brace, amplitude=2mm, mirror}, line width=0.7pt, ink]
     (-15,-38.5) -- (15,-38.5) node[midway, below=3mm, font=\footnotesize, text width=40mm, align=center]
     {\textbf{the biggest gap:} better labels};
\end{tikzpicture}
\caption{Our pipeline and the winners' pipeline, side by side. Both have the same four parts. The largest
difference is in the first one: how good the training labels are.}
\label{fig:story}
\end{figure}

\paragraph{The story in six sentences.}
\begin{enumerate}
\item The task: look at one knee MRI exam and give a probability for each of 12 problems (a torn ligament,
      a torn meniscus, fluid in the joint, \ldots).
\item Only 58 of the 4,407 training exams were labelled by doctors; the rest only have a free-text report.
\item So everybody first had to \textbf{turn reports into labels}, and then train an image model on those labels.
\item We did this with an open language model that read the reports, and trained small image models: 0.918 on
      our own, 0.927 on the final leaderboard with the help of a public notebook.
\item The winners let AI agents read the reports \emph{and look at the images}, and also asked them to draw boxes
      around the anatomy; their models scored 0.960.
\item Most of the gap comes from the labels, a smaller part from telling the model where to look.
\end{enumerate}
